% !TEX program = pdflatex
% ============================================================================
% Isiliel — Programming Manual/Tutorial
% Single file, compilable on Overleaf with pdfLaTeX.
% ============================================================================
\documentclass[11pt,a4paper,oneside]{report}

\usepackage[utf8]{inputenc}
\usepackage[T1]{fontenc}
\usepackage[english]{babel}
\usepackage{lmodern}
\usepackage{microtype}
\usepackage[a4paper,top=2.4cm,bottom=2.4cm,left=2.5cm,right=2.5cm]{geometry}
\usepackage{amsmath,amssymb}
\usepackage{booktabs}
\usepackage{longtable}
\usepackage{array}
\usepackage{enumitem}
\usepackage{xcolor}
\usepackage{listings}
\usepackage{textcomp}
\usepackage{upquote}
\usepackage{fancyhdr}
\usepackage{titlesec}
\usepackage{hyperref}

\definecolor{isilielblue}{RGB}{20,70,120}
\definecolor{isilielgreen}{RGB}{30,110,60}
\definecolor{isilielred}{RGB}{150,40,40}
\definecolor{isilielgray}{RGB}{246,246,248}

\hypersetup{
  colorlinks=true,
  linkcolor=isilielblue,
  urlcolor=isilielblue,
  citecolor=isilielblue,
  pdftitle={Isiliel v1 — Programming Manual/Tutorial},
  pdfauthor={EEPS and MiMo-V2.6},
  pdfsubject={Embedded scripting language}
}

\pagestyle{fancy}
\fancyhf{}
\fancyhead[L]{\small Isiliel v1 --- Programming Manual/Tutorial}
\fancyhead[R]{\small\thepage}
\renewcommand{\headrulewidth}{0.4pt}

\titleformat{\chapter}[display]
  {\normalfont\huge\bfseries\color{isilielblue}}
  {\filright\Large Chapter \thechapter}
  {1ex}
  {\filright}
\titlespacing*{\chapter}{0pt}{-10pt}{24pt}

\setlength{\parindent}{0pt}
\setlength{\parskip}{0.55em}
\setlength{\emergencystretch}{2em}
\setlist[itemize]{leftmargin=1.6em,itemsep=0.15em,topsep=0.25em}
\setlist[enumerate]{leftmargin=1.8em,itemsep=0.15em,topsep=0.25em}

\lstdefinelanguage{Isiliel}{
  morekeywords={if,else,elseif,while,for,return,break,continue,void,int,float,
                fn,class,extends,new,super,include,import,type,global,thread,
                start,stop,pause,resume,repeat,until,goto,nil,true,false,local},
  morecomment=[l]{//},
  morestring=[b]",
  sensitive=true,
  alsoletter={\_}
}

\lstdefinestyle{isiliel}{
  language=Isiliel,
  basicstyle=\ttfamily\small,
  keywordstyle=\color{isilielblue}\bfseries,
  commentstyle=\color{isilielgreen}\itshape,
  stringstyle=\color{isilielred},
  numbers=none,
  frame=single,
  rulecolor=\color{gray!45},
  backgroundcolor=\color{isilielgray},
  breaklines=true,
  breakatwhitespace=false,
  columns=fullflexible,
  keepspaces=true,
  showstringspaces=false,
  tabsize=4,
  xleftmargin=0.8em,
  xrightmargin=0.8em,
  aboveskip=0.7em,
  belowskip=0.7em
}
\lstset{style=isiliel}

\newcommand{\isiliel}{\textsc{Isiliel}}
\newcommand{\code}[1]{\texttt{#1}}
\newcommand{\file}[1]{\texttt{#1}}

\newcolumntype{L}[1]{>{\raggedright\arraybackslash}p{#1}}
\newcolumntype{C}[1]{>{\centering\arraybackslash}p{#1}}

\begin{document}

% ---------------------------------------------------------------------------
% Cover
% ---------------------------------------------------------------------------
\hypersetup{pageanchor=false}
\begin{titlepage}
\centering
\vspace*{1.5cm}
{\Huge\bfseries\color{isilielblue} Isiliel v1\par}
\vspace{0.4cm}
{\LARGE Programming Manual/Tutorial\par}
\vspace{0.8cm}
{\large Embedded scripting language with a stack VM,\\
classes, tables, coroutines, graphics, and a REPL\par}
\vspace{2cm}
\rule{0.75\textwidth}{0.6pt}\par
\vspace{0.8cm}
{\large EEPS and MiMo-V2.6\par}
{\large Version 1.0\par}
{\large \today\par}
\vspace{0.8cm}
\rule{0.75\textwidth}{0.6pt}\par
\vfill
{\small This document describes the language as implemented in version 1.0.\\
Whenever the project documentation and the code diverge, the code is the source of truth.}
\end{titlepage}
\hypersetup{pageanchor=true}

\tableofcontents
\clearpage

\chapter{About This Manual}
\label{cap:sobre}

This is a practical manual and an introductory reference for the
\isiliel{} language. It is written for people who want to \emph{use} the
language: write scripts, organize code into functions and modules, use
tables and classes, draw on screen, control coroutines, and embed the
interpreter in a C++ program.

The text assumes that you know how to program in some imperative language
(C, Python, Lua, or JavaScript). You do not need to know the internal
implementation, but some chapters mention behavioral details to help you
avoid pitfalls.

\section{How to Read}
\begin{itemize}
  \item Chapters \ref{cap:inicio} through \ref{cap:funcoes} form the
        basic tutorial: installation, values, variables, control flow,
        and functions.
  \item Chapters \ref{cap:strings} through \ref{cap:metatables} cover the
        data types and object-oriented programming.
  \item Chapters \ref{cap:modulos} through \ref{cap:embed} cover
        modules, errors, concurrency, the standard library, graphics, and
        embedding.
  \item The appendices contain the grammar summary, a quick builtin
        reference, and complete programs.
\end{itemize}

\section{Conventions}
\begin{itemize}
  \item Gray-background blocks are \isiliel{} code. Comments use
        \code{//}, and example text avoids accents inside code so it
        remains compatible with any editor.
  \item \code{println}, \code{printf}, and function names appear in a
        monospaced font.
  \item Paths such as \file{scripts/script\_1.isi} are relative to the
        project root.
  \item ``v1.0'' refers to the version described in this manual.
\end{itemize}

\section{Note on Known Differences}
The language is evolving. Some points in the project documentation
(\file{docs/roadmap.md} and \file{docs/linguagem.md}) describe the
\emph{planned} behavior; this manual describes the \emph{current}
behavior. The most important differences are collected in
Chapter~\ref{cap:limitacoes}. In particular, in v1.0:
\begin{itemize}
  \item \code{true} and \code{false} are numeric literals (\code{1} and
        \code{0}); there is no distinct runtime boolean type.
  \item \code{\&\&} and \code{||} evaluate both operands and produce
        \code{1}/\code{0}; there is no short-circuiting.
  \item \code{!} is bitwise negation, not logical negation.
  \item \code{pairs}/\code{ipairs} return lists; there is no
        \code{for k,v in ...}.
\end{itemize}

% ===========================================================================
\part{Getting Started}
% ===========================================================================

\chapter{Overview}
\label{cap:visao}

\isiliel{} is a dynamic scripting language implemented in C++17 with the
pipeline:

\begin{center}
\fbox{\parbox{0.9\textwidth}{\centering
source \ $\rightarrow$ \ lexer \ $\rightarrow$ \ bytecode compiler
\ $\rightarrow$ \ stack VM}}
\end{center}

The same executable provides a file \emph{runner}, an interactive REPL,
a disassembler, and an embedding API. The language was designed to be
small and portable, with a focus on embedded systems, but it includes
features found in general-purpose scripting languages.

\section{Main Features}
\begin{itemize}
  \item Exact 64-bit integers and double-precision floating-point
        numbers; complex numbers with the \code{j} suffix.
  \item Immutable strings with escapes, indexing, slicing, and patterns.
  \item Arduino String API (free functions) and C99 \file{math.h}
        functions/constants.
  \item 0-based lists/matrices with slicing, \code{append}, \code{sort},
        and basic linear algebra.
  \item Tables (dictionaries) with text, numeric, or boolean keys.
  \item First-class functions, multiple returns, varargs, and anonymous
        closures.
  \item Structs and classes with single inheritance, an \code{init}
        constructor, and \code{super}.
  \item Metatables and metamethods (\code{\_\_index}, \code{\_\_call},
        \code{\_\_add}, \code{\_\_eq}, \code{\_\_tostring}, \code{\_\_gc}).
  \item Cycle collection by \emph{trial deletion}.
  \item Modules through \code{include}, \code{require}, \code{load}, and
        \code{dofile}; native \code{.mod} modules.
  \item System threads and cooperative coroutines.
  \item Processing-style 2D graphics with a headless software backend or
        raylib.
  \item REPL, automated tests, ASan/UBSan, fuzzer, and benchmarks.
\end{itemize}

\section{Execution Model}
A script is compiled entirely to bytecode before execution. The VM is
stack-based and keeps global and local symbols in a symbol table. The
compiler resolves names at compile time when possible; direct calls use
the function address, while function values use indirect calls.

Running a file prints diagnostic banners and returns exit code \code{0}
on success or \code{1} on error. The REPL suppresses the banners and
keeps symbols and bytecode between lines.

\section{Target Audience and Portability}
The main implementation targets desktop systems, but the code avoids
mandatory operating-system dependencies. The graphics backend is
abstract (\code{software} or \code{raylib}), and the \code{make embedded}
profile uses \code{-Os} without raylib. Threads and files are optional
in embedded environments; see \file{docs/proposta\_embarcada.md} for the
design decisions.

% ---------------------------------------------------------------------------
\chapter{Installation and First Steps}
\label{cap:inicio}

\section{Building}
From the project root:

\begin{lstlisting}[language=bash]
make                 # build/isirun in debug mode (-O0 -g3)
make test            # full test suite (164 cases in v1.0)
make release         # -O2 -DNDEBUG -flto
make embedded        # -Os, no raylib (embedded profile)
make asan            # suite with AddressSanitizer + UBSan
make fuzz            # deterministic fuzzer
make bench           # benchmark against Lua 5.4
\end{lstlisting}

The \code{raylib} backend is detected through \code{pkg-config}; use
\code{make USE\_RAYLIB=1} to force it or \code{make USE\_RAYLIB=0} to
disable it. Without raylib, everything runs on the \code{software}
backend (headless, writing \code{.ppm} images).

\section{Running a Script}
\begin{lstlisting}[language=bash]
./build/isirun scripts/script_1_matriz_lu.isi
./build/isirun --disasm scripts/script_1_matriz_lu.isi   # shows the bytecode
./build/isirun                                 # interactive REPL
\end{lstlisting}

A script is a text file (usually \code{.isi}) with statements terminated
by \code{;}. First example:

\begin{lstlisting}
// hello.isi
nome = "mundo";
println("Ola, " + nome + "!");
\end{lstlisting}

Output:
\begin{lstlisting}[language=bash]
Ola, mundo!
\end{lstlisting}

\section{The REPL}
With no arguments, the executable starts a REPL. State (variables,
functions, and bytecode) persists between lines. Blocks are entered over
multiple lines; the \code{...} prompt continues until the braces are
closed. A simple statement ends with \code{;}; a block ends with
\code{\}}. \code{exit}, \code{quit}, or EOF terminate the session.

\begin{lstlisting}[language=bash]
>>> a = 2 + 3;
>>> println(a);
5
>>> f(x) { return x * x; }
>>> println(f(7));
49
\end{lstlisting}

Compilation or execution errors are reported but do not end the session.
The REPL enforces growth limits to prevent unlimited memory consumption
(default: 250000 instructions and 50000 symbols). Adjust them with the
environment variables \code{MYINTERP\_REPL\_CODE\_LIMIT} and
\code{MYINTERP\_REPL\_SYMBOLS\_LIMIT}; \code{0} removes the limit.

\section{First Complete Program}
\begin{lstlisting}
// media.isi: reads two numbers and computes the average
a = 7;
b = 9;
media = (a + b) / 2;
printf("media = %g\n", media);

if (media >= 7) {
    println("aprovado");
} else {
    println("reprovado");
}
\end{lstlisting}

Note that \code{/} always produces a floating-point number. For integer
division use \code{floordiv(a, b)}.

\section{Errors and Traceback}
Compilation errors include the line:
\begin{lstlisting}[language=bash]
ERRO COMPILACAO: Esperado ';' apos a expressao (linha 4)
\end{lstlisting}

Execution errors include the line and, when available, a traceback:
\begin{lstlisting}[language=bash]
ERRO EM TEMPO DE EXECUCAO: Divisao por zero (linha 7)
  em quebra (linha 9)
\end{lstlisting}

Use \code{pcall} to catch runtime errors (Section~\ref{sec:pcall}).

% ---------------------------------------------------------------------------
\chapter{Lexical Structure}
\label{cap:lexico}

\section{Comments and Statements}
Comments start with \code{//} and run to the end of the line. There are no
block comments. Simple statements end with \code{;}. Some statements
(function definitions, \code{include}, \code{import},
\code{start}/\code{stop}/\code{pause}/\code{resume}) accept an optional
semicolon; use \code{;} for consistency.

\section{Identifiers and Keywords}
Identifiers start with a letter or \code{\_} and may contain letters,
digits, and \code{\_}. The reserved words are:

\begin{center}
\begin{tabular}{llll}
\toprule
\code{if} & \code{else} & \code{elseif} & \code{while} \\
\code{for} & \code{return} & \code{break} & \code{continue} \\
\code{void} & \code{int} & \code{float} & \code{fn} \\
\code{class} & \code{extends} & \code{new} & \code{super} \\
\code{include} & \code{import} & \code{type} & \code{global} \\
\code{thread} & \code{start} & \code{stop} & \code{pause} \\
\code{resume} & \code{repeat} & \code{until} & \code{goto} \\
\code{nil} & & & \\
\bottomrule
\end{tabular}
\end{center}

\code{true}, \code{false}, \code{i32}, \code{f64}, \code{bool}, and
\code{str} are not keywords: they are literals or prefixes recognized by
the parser. That is why \code{type(x)} (the function) and
\code{type Nome \{\}} (the struct definition) coexist without conflict.

\section{Numeric Literals}
\begin{itemize}
  \item Decimal integers: \code{0}, \code{42}, \code{9223372036854775807}.
        Literals that fit in \code{int64} are exact.
  \item Floating point: \code{3.14}, \code{0.5}. The lexer does not accept
        scientific notation: \code{1e9} is read as \code{1} followed by the
        identifier \code{e9}.
  \item Complex numbers: suffix \code{j} or \code{J}, as in Python:
        \code{3j}, \code{1+2j}, \code{2.5J}.
\end{itemize}

\section{Strings}
Strings are delimited by double quotes. Accepted escapes:
\code{\textbackslash n}, \code{\textbackslash t},
\code{\textbackslash r}, \code{\textbackslash\textbackslash},
\code{\textbackslash "}, and \code{\textbackslash 0}. Unknown escapes
keep the backslash (Python style): \code{"\textbackslash q"} becomes
\code{\textbackslash q}. Strings are immutable.

\begin{lstlisting}
s = "linha 1\nlinha 2\tfim";
println(len(s));      // number of characters
println(s[0]);        // first character
println(s[-1]);       // last character
\end{lstlisting}

\section{Operators}
\begin{center}
\begin{tabular}{ll}
\toprule
Category & Operators \\
\midrule
Arithmetic & \code{+  -  *  /  \%} \\
Bitwise & \code{\&  |  \^{}  \textasciitilde  <<  >>} \\
Comparison & \code{==  !=  <  >  <=  >=} \\
Logical & \code{\&\&  ||} \\
Assignment & \code{=  +=  -=  *=  /=} \\
Access & \code{.  []  ()} \\
\bottomrule
\end{tabular}
\end{center}

There is no \code{and}, \code{or}, \code{not}, \code{\#}, \code{..},
or \code{//} operator. \code{//} is a comment; for integer division use
\code{floordiv(a,b)}. For string concatenation use \code{+}. For logical
negation use an explicit comparison, such as \code{x == 0} or
\code{x == nil}.

% ===========================================================================
\part{The Language}
% ===========================================================================

\chapter{Values, Types, and Expressions}
\label{cap:valores}

\section{Available Types}
The table below shows the values the programmer sees and what
\code{type(x)} returns.

\begin{center}
\begin{longtable}{L{3.4cm}L{7.4cm}L{3.2cm}}
\toprule
Type & Example & \code{type(x)} \\
\midrule
Integer & \code{30}, \code{-7}, \code{0x1F} & \code{number} \\
Floating point & \code{3.14}, \code{7 / 2} & \code{number} \\
String & \code{"abc"} & \code{string} \\
Boolean & \code{true}, \code{false} (1/0 in v1.0) & \code{number} \\
List/matrix & \code{[1, 2, 3]}, \code{[[1,2],[3,4]]} & \code{array} \\
Complex & \code{3j}, \code{1+2j} & \code{complex} \\
Struct/class & \code{Ponto\{1, 2\}} & \code{struct} \\
Table & \code{\{nome = "Ana"\}} & \code{table} \\
Function & \code{fn(x) \{ return x; \}} & \code{function} \\
Coroutine & \code{cocreate(f)} & \code{thread} \\
Absent & \code{nil} & \code{nil} \\
\bottomrule
\end{longtable}
\end{center}

\section{Integers and Floating Point}
Integers that fit in \code{int64} are exact. \code{+}, \code{-}, and
\code{*} between integers wrap around at 64 bits. \code{/} always
returns a \code{double}:

\begin{lstlisting}
println(7 / 2);       // 3.5
println(7 % 2);       // 1
println(floordiv(7, 2)); // 3
println(9223372036854775807 + 1); // -9223372036854775808
\end{lstlisting}

Mixing an integer and a \code{double} promotes to \code{double}. Bitwise
operators require integers; \code{5.5 \& 1} is an error.

\section{Strings}
Strings are immutable. The \code{+} operator concatenates two strings;
mixing a string and a number is an error (there is no silent coercion).
Use \code{str(x)} or \code{printf} to convert.

\begin{lstlisting}
nome = "Ana";
println("Ola, " + nome);     // Hello, Ana
println("abc" == "abc");     // 1
println("abc" < "abd");      // 1 (lexicographic)
\end{lstlisting}

\section{Truth and Falsity}
The table below describes the current implementation.

\begin{center}
\begin{tabular}{L{7cm}L{7cm}}
\toprule
Value & Truthy? \\
\midrule
\code{nil} & no \\
\code{0}, \code{0.0}, \code{0j} & no \\
\code{false} (which is \code{0}) & no \\
Any string, including \code{""} & yes \\
List, table, struct, function, coroutine & yes \\
Nonzero number & yes \\
\bottomrule
\end{tabular}
\end{center}

The project documentation lists the empty string as false, but the
current code treats every string as true. Use \code{len(s) == 0} if you
want to test for an empty string.

\section{Logical Operators}
In v1.0, \code{\&\&} and \code{||} evaluate both operands and return
\code{1} or \code{0}. There is no short-circuiting. This means that
\code{f() \&\& g()} always calls both functions. For conditional side
effects, use \code{if}:

\begin{lstlisting}
if (a != 0) { g(); }
\end{lstlisting}

\code{!} and \code{\textasciitilde} are the same bitwise negation
operator. Do not use \code{!x} as logical negation: \code{!0} is
\code{-1}, which is true. To negate a condition, write \code{x == 0} or
\code{x == nil}; \code{if (!(x))} is not recommended.

\section{Precedence}
From lowest to highest precedence:

\begin{center}
\begin{tabular}{ll}
\toprule
Level & Operators \\
\midrule
1 & assignment (\code{= += -= *= /=}) \\
2 & \code{||} \\
3 & \code{\&\&} \\
4 & \code{|} \\
5 & \code{\^{}} \\
6 & \code{\&} \\
7 & \code{== !=} \\
8 & \code{< > <= >=} \\
9 & \code{<< >>} \\
10 & \code{+ -} \\
11 & \code{* / \%} \\
12 & unary \code{-}, \code{!}, \code{\textasciitilde} \\
13 & postfix \code{() [] .} \\
\bottomrule
\end{tabular}
\end{center}

Use parentheses whenever in doubt.

\section{Assignment and Conversions}
Assignment can be used as an expression. The existing compound operators
are \code{+=}, \code{-=}, \code{*=}, and \code{/=}; there is no
\code{\%=}. Useful conversions:

\begin{lstlisting}
println(tonumber("42"));      // 42
println(tonumber("3.5"));     // 3.5
println(tonumber("ff", 16));  // 255
println(str(42));             // "42"
println(tostring(nil));       // nil
println(type([1,2]));         // array
\end{lstlisting}

An invalid \code{tonumber} returns \code{nil}. \code{str}/\code{tostring}
produce the textual representation; \code{println} of a string does not
add quotes, and neither does \code{str} of a string.

% ---------------------------------------------------------------------------
\chapter{Variables and Scope}
\label{cap:variaveis}

\section{Globals and Implicit Locals}
At the top level of a script, an assignment creates or changes a global
variable:

\begin{lstlisting}
x = 10;
println(x);       // 10
\end{lstlisting}

Inside a function, an assignment to a name that is not a parameter
creates a \emph{local} variable of that function. Reads look for the
local first; if it does not exist, they look for the global. To write to a
global from inside a function, declare it with \code{global}:

\begin{lstlisting}
global contador;
contador = 0;

incrementa() {
    contador += 1;
}

incrementa();
println(contador);   // 1
\end{lstlisting}

Without \code{global}, \code{contador += 1} would create a new local on
each call. This is essential in the \code{setup()} and \code{draw()}
callbacks.

\section{Optional Scalar Declarations}
The language is dynamic by default. Optionally, you can annotate a
variable's scalar type in C/Rust style. The annotation applies coercion on
every write:

\begin{lstlisting}
i32 a = 42;
u8 porto = 255;
f32 p = 3.14;
f64 q = 3.14;
bool ligado = 1;
str nome = "Ana";
int i = 7;          // alias of i32
float f = 2.5;      // alias of f64
String s = "abc";   // alias of str
\end{lstlisting}

\begin{center}
\begin{tabular}{L{4.5cm}L{9.5cm}}
\toprule
Type & Coercion rule \\
\midrule
\code{i8..i64}, \code{isize} & truncates toward zero and checks range \\
\code{u8..u64}, \code{usize} & truncates and requires non-negative \\
\code{f32} & rounds to float32 precision \\
\code{f64} & normal double \\
\code{bool} & nonzero becomes 1; zero becomes 0 \\
\code{str} & requires a string \\
\bottomrule
\end{tabular}
\end{center}

The type prefix on functions, parameters, and fields is accepted and
\emph{discarded}: the language remains dynamic at those points. Only
scalar variables declared with a type have automatic coercion.

\section{Multiple Assignment}
You can assign several values at once. All values are evaluated before any
write:

\begin{lstlisting}
a = 0; b = 0;
a, b = 3, 4;       // a=3, b=4
a, b = b, a;       // true swap
p, q, r = 1, 2, 3;
\end{lstlisting}

The number of targets must equal the number of values, except when a
single call provides several values:
\begin{lstlisting}
fn par() { return 10, 20; }
x, y = par();       // x=10, y=20
\end{lstlisting}

Targets are simple variables; do not use \code{a.b, c = ...}.

\section{Recursion and Local Cleanup}
Each function call saves and clears its locals, allowing recursion and
preventing state from leaking between calls:

\begin{lstlisting}
fatorial(n) {
    if (n <= 1) { return 1; }
    return n * fatorial(n - 1);
}
println(fatorial(10));   // 3628800
\end{lstlisting}

\section{REPL}
In the REPL, globals and functions persist between lines. A function's
locals do not leak to the top level. Redefining a function or variable is
allowed; for \code{type}, redefinition clears the type's fields and
methods.

% ---------------------------------------------------------------------------
\chapter{Control Flow}
\label{cap:controle}

\section{if / elseif / else}
\begin{lstlisting}
n = 7;
if (n % 2 == 0) {
    println("par");
} elseif (n % 3 == 0) {
    println("multiplo de 3");
} else {
    println("outro");
}
\end{lstlisting}

\code{elseif} is a single word; \code{else if} also works. Conditions
require parentheses.

\section{while}
\begin{lstlisting}
i = 0;
while (i < 3) {
    println(i);
    i += 1;
}
\end{lstlisting}

\section{Classic for}
The only \code{for} in v1.0 is the C-style loop:

\begin{lstlisting}
for (i = 0; i < 5; i += 1) {
    println(i);
}
\end{lstlisting}

All three clauses are optional: \code{for (;;) \{ ... \}} is an infinite
loop. There is no \code{for k,v in pairs(t)}; to iterate over tables, use
\code{pairs(t)} and a \code{for} over the resulting list
(Chapter~\ref{cap:tabelas}).

\section{repeat / until}
The body executes at least once; the condition is tested at the end:

\begin{lstlisting}
i = 3;
repeat {
    println(i);
    i -= 1;
} until (i <= 0);
\end{lstlisting}

\code{until} accepts the condition with or without parentheses.

\section{break, continue, and goto}
\code{break} exits the innermost loop. \code{continue} goes to the next
iteration; in a \code{for} loop, the increment runs before the condition.

\begin{lstlisting}
for (i = 0; i < 10; i += 1) {
    if (i == 3) { continue; }
    if (i == 7) { break; }
    println(i);
}
\end{lstlisting}

Lua-style \code{goto} and labels work within the same function:

\begin{lstlisting}
i = 0;
::inicio::
i += 1;
if (i < 3) { goto inicio; }
println(i);   // 3
\end{lstlisting}

The label must exist in the same function; duplicates are a compilation
error.

\section{Tail Calls}
\code{return f(args);} with a direct call to a script function is a
\emph{proper tail call}: the current frame is reused. This allows
constant-space tail recursion. Indirect calls and builtins do not use this
optimization.

% ---------------------------------------------------------------------------
\chapter{Functions}
\label{cap:funcoes}

\section{Definition and Call}
A function can be defined with or without the \code{fn} prefix. Type
prefixes (\code{void}, \code{int}, \code{i32}, \code{float}, etc.) are
accepted and ignored.

\begin{lstlisting}
soma(a, b) {
    return a + b;
}

fn dobro(x) { return x * 2; }
int triplo(int x) { return x * 3; }

println(soma(2, 3));   // 5
println(dobro(4));     // 8
println(triplo(3));    // 9
\end{lstlisting}

Functions are global. Parameters and variables assigned in the body are
local. A function without \code{return} returns zero values; in a
single-value context, the result is \code{nil}.

\section{Multiple Returns}
\begin{lstlisting}
divmod(a, b) {
    return a / b, a % b;
}

q, r = divmod(7, 2);
println(q);   // 3.5
println(r);   // 1
\end{lstlisting}

In a single-value context, only the first return value is used. In
multiple assignment, values are distributed in order; missing values
become \code{nil}, and extras are discarded. \code{return f();} forwards
all values from the callee (tail call).

\section{Varargs}
The last parameter can be \code{...}. Inside the function, \code{...}
represents the extra arguments.

\begin{lstlisting}
soma_var(first, ...) {
    total = first;
    for (i = 0; i < select("#", ...); i += 1) {
        total += select(i + 1, ...);
    }
    return total;
}

println(soma_var(1, 2, 3, 4));   // 10
\end{lstlisting}

\code{select("\#", ...)} returns the number of extras;
\code{select(n, ...)} returns from the n-th onward (1-based; negative
counts from the end). \code{unpack(lista[, i[, j]])} returns the elements
(0-based indices). \code{...} can only be the last argument of a call or
of a \code{return}; in that position it expands to all values. Otherwise,
\code{...} in an ordinary expression is the first extra value.

\section{Functions as Values}
\begin{lstlisting}
dobro(n) { return n * 2; }
f = dobro;
println(f(4));        // 8
println(f == dobro);  // 1 (same function)
println(f);           // <function dobro>
\end{lstlisting}

Calling a variable that holds a function uses an indirect call.
Functions cannot be used in arithmetic or converted to scalar types.

\section{Anonymous Closures}
An anonymous function \code{fn(...) \{ ... \}} captures local variables
from the enclosing scope by reference:

\begin{lstlisting}
contador() {
    n = 0;
    return fn() {
        n += 1;
        return n;
    };
}

c1 = contador();
c2 = contador();
println(c1());   // 1
println(c1());   // 2
println(c2());   // 1
\end{lstlisting}

Two closures created in the same frame share the same upvalue.

\paragraph{v1.0 limitations.}
Only anonymous functions capture. Named functions and methods do not
capture outer locals; use an inner anonymous function. Capture is lexical
but not forward-looking: declare and assign the local before creating the
closure. There is no IIFE: \code{fn(x)\{...\}(1)} is not supported;
assign it to a variable first. Scalar coercion of a captured local is not
reapplied on reads/writes through the upvalue.

\section{global inside Functions}
As seen in Chapter~\ref{cap:variaveis}, use \code{global nome;} to share
state between functions and callbacks. Without it, assignment creates a
local.

\section{pcall and Traceback}
\label{sec:pcall}
\code{pcall(f, ...)} runs \code{f(...)} and catches runtime errors. It
returns a 1x2 matrix: \code{[ok, value]}.

\begin{lstlisting}
fn quebra() { return 1 / 0; }

r = pcall(quebra);
println(r[0]);   // 0 (failure)
println(r[1]);   // message + traceback

r = pcall(fn(a, b) { return a + b; }, 2, 3);
println(r[0]);   // 1
println(r[1]);   // 5
\end{lstlisting}

If the first argument is not a function, the result is
\code{[0, "pcall: valor nao e uma funcao"]}. \code{pcall} is a special
call form, not a value: do not write \code{g = pcall}. Unhandled errors
show the line and frames in the REPL/runner.

\section{Builtins as Functions}
Builtins such as \code{println}, \code{len}, \code{str}, \code{abs},
and \code{sqrt} are called by name. They are not first-class values in
v1.0; to pass an operation as an argument, wrap it in an anonymous
function:

\begin{lstlisting}
aplica(f, x) { return f(x); }
println(aplica(fn(v) { return v * v; }, 5));   // 25
\end{lstlisting}

% ---------------------------------------------------------------------------
\chapter{Strings}
\label{cap:strings}

\section{Basic Operations}
\begin{lstlisting}
s = "Ola, mundo";
println(len(s));       // 11
println(s[0]);         // O
println(s[-1]);        // o
println(upper(s));     // HELLO, WORLD
println(lower(s));     // hello, world
println(trim("  x  ")); // x
println(rep("ab", 3)); // ababab
println(reverse(s));   // reversed text
\end{lstlisting}

Strings are immutable: \code{s[0] = "X"} is an error. Negative indices
count from the end. \code{len} counts characters (bytes in internal
storage; for UTF-8 code points use \code{utf8\_len}).

\section{Slices}
Slices work on strings and lists. The start is inclusive and the end is
\emph{exclusive}; bounds are clamped and do not raise an error.

\begin{lstlisting}
s = "abcdef";
println(s[1:4]);   // bcd
println(s[:3]);    // abc
println(s[3:]);    // def
println(s[-2:]);   // ef
\end{lstlisting}

\code{sub(s, i[, j])} is different: it uses 0-based indices and the end
is \emph{inclusive}. See Appendix~\ref{ap:builtins}.

\section{Comparison and Concatenation}
\code{+} concatenates two strings. \code{==} and \code{!=} compare
content; \code{< > <= >=} are lexicographic between two strings.
Comparing a string with a number is an error.

\begin{lstlisting}
println("abc" + "def");   // abcdef
println("abc" == "abc");  // 1
println("abc" != "abd");  // 1
\end{lstlisting}

\section{Conversion and Formatting}
\code{str(x)} and \code{tostring(x)} convert to text.
\code{format(fmt, ...)} returns a formatted string; \code{printf} prints
it. The accepted conversions are \code{\%s}, \code{\%d}, \code{\%i},
\code{\%f}, \code{\%g}, \code{\%j}, \code{\%c}, and \code{\%\%}.
There is no width, precision, or flags: \code{\%.2f} is an error.

\begin{lstlisting}
nome = "Ana";
idade = 30;
msg = format("%s tem %d anos", nome, idade);
println(msg);              // Ana tem 30 anos
printf("[%s]\n", "ok");    // [ok]
\end{lstlisting}

\section{Arduino String API}
The functions below are free functions (not methods) and preserve string
immutability; they all accept text and return text or a number depending
on the operation.

\begin{lstlisting}
s = "  Arduino  ";
println(trim(s));                    // Arduino
println(String(3.14159, 2));         // 3.14
println(String(255, 16));            // ff
println(charAt("abc", 1));           // b
println(indexOf("arduino", "du"));   // 2
println(toUpperCase("abc"));         // ABC

buf = [0, 0, 0, 0];
n = toCharArray("Az", buf);          // 2; buf = [65, 122, 0, 0]
println(buf[0]);                     // 65
\end{lstlisting}

Without a buffer, \code{getBytes}/\code{toCharArray} return a 1D list of
bytes; with a buffer, they return the number of bytes copied.
\code{getBytes} does not write a NUL; \code{toCharArray} writes a NUL if
space remains. \code{String(float)} uses 2 decimal places;
\code{String(int, base)} accepts base 2..36. There are also
\code{compareTo}, \code{c\_str}, \code{endsWith}, \code{equals},
\code{equalsIgnoreCase}, \code{lastIndexOf}, \code{length}, \code{remove},
\code{replace}, \code{reserve}, \code{setCharAt}, \code{startsWith},
\code{substring}, \code{toDouble}, \code{toFloat}, \code{toInt},
\code{toLowerCase}, and \code{toUpperCase}; \code{concat(s, ...)}
concatenates text when the first argument is a string.

\section{Patterns}
The language implements a subset of Lua patterns. Indices are 0-based; a
failed \code{match} returns \code{""} (empty string) for compatibility.

\begin{lstlisting}
println(match("hello 123", "%d+"));       // 123
println(match("abc123", "(%a+)(%d+)"));   // abc (first capture)
println(gsub("a1b2", "%d", "#"));         // a#b#
println(gsub("abc", "(%a)", "[%1]"));     // [a][b][c]

lista = gmatch("a1 b22", "%d+");          // list ["1", "22"]
for (i = 0; i < len(lista); i += 1) {
    println(lista[i]);
}
\end{lstlisting}

Supported classes: \code{. \%a \%c \%d \%g \%l \%p \%s \%u \%w \%x \%z}
(uppercase inverts), sets \code{[...]} with ranges and \code{\^{}},
quantifiers \code{* + - ?}, anchors \code{\^{} \$}, captures
\code{(...)}, \code{\%bxy}, \code{\%f[set]}, and backreferences
\code{\%1..\%9}. \code{gsub} accepts textual replacement only.

\section{UTF-8}
\begin{lstlisting}
s = "cafe";
println(utf8_len("caf\u00e9"));       // 4 code points (example)
println(utf8_codepoint("A", 0));      // 65
println(utf8_char(65));               // A
\end{lstlisting}

\code{utf8\_len(s)} counts code points; \code{utf8\_codepoint(s[, i])}
returns the code point at byte \code{i} or \code{-1};
\code{utf8\_char(cp)} encodes a valid code point. The source above uses
\code{\textbackslash u} only as an illustration; the v1.0 lexer does not
interpret \code{\textbackslash u}, so prefer literal UTF-8 strings.

% ---------------------------------------------------------------------------
\chapter{Lists and Matrices}
\label{cap:listas}

\section{1D Lists}
A list is a 1xN matrix, indexed from 0.

\begin{lstlisting}
v = [10, 20, 30];
println(v[0]);        // 10
println(v[-1]);       // 30
println(len(v));      // 3

v[1] = 99;
append(v, 40);        // mutates and returns the list
println(v);           // 1x4 matrix: 10, 99, 30, 40
\end{lstlisting}

\code{append} only works on 1D lists. \code{sort(v)} sorts numbers in
ascending order (in place); \code{concat(v[, sep])} joins the elements
into text.

\section{Slices and Copies}
Slices return a \emph{copy} (mutating the copy does not change the
original). The end is exclusive.

\begin{lstlisting}
v = [0, 1, 2, 3, 4];
w = v[1:4];       // [1, 2, 3]
println(len(w));  // 3
w[0] = 99;
println(v[1]);    // 1 (original intact)
\end{lstlisting}

Assigning to a slice (\code{v[1:3] = ...}) is not supported.

\section{2D Matrices}
Lists of lists form matrices. \code{mat\_index(M, i, j)} uses 1-based
indices (the linear algebra convention), while direct indexing
\code{M[i][j]} uses 0-based indices.

\begin{lstlisting}
M = [[1, 2], [3, 4]];
println(M[0][1]);          // 2
println(mat_index(M, 1, 2)); // 2
println(rows(M));          // 2
println(cols(M));          // 2
\end{lstlisting}

Available operations: \code{det}, \code{trace}, \code{inv},
\code{transpose}, \code{dot}, \code{cross}, \code{matmul},
\code{identity}, \code{zeros}, \code{eigvals}, and \code{eigvecs}. The
operators \code{+ - * /} and \code{@} (scalar) work on matrices.

% ---------------------------------------------------------------------------
\chapter{Tables}
\label{cap:tabelas}

\section{Creating and Accessing}
Tables are key-value maps with reference semantics. The literal uses
\code{key = value} pairs; literal keys are identifiers.

\begin{lstlisting}
pessoa = {nome = "Ana", idade = 30};
println(pessoa.nome);          // Ana
println(pessoa["idade"]);      // 30
pessoa.idade += 1;
pessoa["cidade"] = "Recife";   // dynamic key
println(pessoa.cidade);        // Recife
\end{lstlisting}

Valid keys are text, number, or boolean. A missing key returns
\code{nil}; assigning \code{nil} deletes the key.

\begin{lstlisting}
println(pessoa.cargo);         // nil
println(has(pessoa, "cargo")); // 0
pessoa.cidade = nil;           // deletes
println(has(pessoa, "cidade"));// 0
\end{lstlisting}

\section{Equality and Reference}
Two tables are equal only if they are the same table (identity):
\begin{lstlisting}
a = {x = 1};
b = a;
c = {x = 1};
println(a == b);   // 1
println(a == c);   // 0
\end{lstlisting}

\section{Iteration}
In v1.0, \code{pairs(t)} returns a list of pairs \code{[k, v]};
\code{ipairs(t)} returns the values of \code{t[0]}, \code{t[1]}, ...
until the first missing key; \code{next(t[, k])} returns the next pair or
an empty list at the end. The order is canonical/deterministic.

\begin{lstlisting}
t = {a = 1, b = 2};
ps = pairs(t);
for (i = 0; i < len(ps); i += 1) {
    par = ps[i];
    println(par[0], "=", par[1]);
}
\end{lstlisting}

You can also iterate over sorted keys:
\begin{lstlisting}
ks = keys(t);
for (i = 0; i < len(ks); i += 1) {
    println(ks[i], t[ks[i]]);
}
\end{lstlisting}

\section{Textual Representation}
\code{println(t)} shows \code{\{key: value\}} with sorted keys. Text keys
appear in quotes; numbers and booleans do not.

\section{Metatables}
Tables and structs can have a metatable; see
Chapter~\ref{cap:metatables}. Table operations also accept
\code{rawget}, \code{rawset}, \code{has}, \code{remove}, \code{keys},
\code{len}, and \code{collectgarbage}.

% ---------------------------------------------------------------------------
\chapter{Structs and Classes}
\label{cap:structs}

\section{Structs}
A struct groups fields and methods. The definition ends with \code{;}.

\begin{lstlisting}
type Ponto {
    x;
    y;

    void mover(dx, dy) {
        x += dx;
        y += dy;
    }
};

p = Ponto{1, 2};
p.mover(3, 4);
println(p.x);   // 4
println(p.y);   // 6
println(p);     // Ponto(x=4, y=6)
\end{lstlisting}

Inside a method, fields are accessed directly by name (\code{x},
\code{y}), without \code{this.}. If you need a local variable with the
same name as a field, choose another name.

The literal \code{Tipo\{v1, v2, ...\}} requires the exact field count, in
definition order (root to leaf). You can also declare
\code{Ponto p = \{1, 2\};}.

\section{Classes and Inheritance}
\code{class} is a struct with single inheritance, a constructor, and
\code{super}.

\begin{lstlisting}
class Conta {
    saldo;
    init(s) { saldo = s; }
    depositar(v) { saldo += v; return saldo; }
}

class Poupanca extends Conta {
    bonus;
    init(s, b) {
        super.init(s);
        bonus = b;
    }
    depositar(v) {
        saldo += v + bonus;
        return saldo;
    }
}

c = new Conta(100);
p = new Poupanca(100, 5);
println(c.depositar(10));   // 110
println(p.depositar(10));   // 115
\end{lstlisting}

\code{new T(args)} requires an \code{init} method (own or inherited).
Fields are flattened in root$\rightarrow$leaf order; methods resolve from
child to parent. \code{super.m(args)} starts at the parent and uses the
same object.

\section{Reference and Equality}
Instances are shared by reference: \code{b = a} points to the same
object. The current implementation compares structs/classes
\emph{structurally} (type and fields); do not rely on identity to
distinguish two instances with the same values. For explicit identity,
add an \code{id} field and compare it.

\section{Structs versus Tables}
Use structs when the set of fields is fixed and part of the type; use
tables for dynamic maps. Structs have methods and inheritance; tables have
dynamic keys and more general metatables.

% ---------------------------------------------------------------------------
\chapter{Metatables and Garbage Collection}
\label{cap:metatables}

\section{setmetatable and getmetatable}
\code{setmetatable(t, mt)} associates a metatable with a table or struct;
\code{nil} removes it. \code{getmetatable(t)} returns the metatable or
\code{nil}.

\section{Supported Metamethods}
\begin{center}
\begin{tabular}{L{3.2cm}L{10.5cm}}
\toprule
Metamethod & Effect \\
\midrule
\code{\_\_index} & lookup when the key/field is missing; may be a table or function \\
\code{\_\_call} & makes the object callable; the object becomes the first argument \\
\code{\_\_add} & covers \code{a + b}, \code{a + k}, \code{x = x + b}, and \code{x = x + k} \\
\code{\_\_eq}, \code{\_\_ne} & equality when both operands have the same type \\
\code{\_\_tostring} & representation used by \code{println}, \code{str}, \code{printf \%s}, \code{error}/\code{assert} \\
\code{\_\_gc} & best-effort finalizer at script end or in \code{collectgarbage()} \\
\bottomrule
\end{tabular}
\end{center}

Examples:
\begin{lstlisting}
base = {saudacao = "ola"};
mt = {__index = base};
obj = setmetatable({nome = "Ana"}, mt);
println(obj.saudacao);        // hello

mtCall = {__call = fn(self, x) { return self.base + x; }};
soma = setmetatable({base = 10}, mtCall);
println(soma(5));             // 15

mtAdd = {__add = fn(a, b) { return a.x + b.x; }};
a = setmetatable({x = 2}, mtAdd);
b = setmetatable({x = 3}, mtAdd);
println(a + b);               // 5

mtStr = {__tostring = fn(self) { return "Pessoa(" + self.nome + ")"; }};
p = setmetatable({nome = "Ana"}, mtStr);
println(p);                   // Pessoa(Ana)
\end{lstlisting}

\section{Cycle Collection}
Tables, lists, structs, functions, and closure environments use reference
counting and a cycle collector based on \emph{trial deletion}. The
collector runs at the end of each script/REPL entry and on demand:

\begin{lstlisting}
t = {};
t.self = t;
t = nil;
println(collectgarbage());   // >= 1

vivo = {};
vivo.self = vivo;
println(collectgarbage());   // 0 (still reachable)
\end{lstlisting}

The collector is not concurrent. Cycles kept alive by stacks of paused
threads/coroutines are treated as roots and may not be collected. The
bytecode of chunks loaded by \code{load}/\code{dofile} accumulates until
the end of execution.

\section{Limitations}
There is still no \code{\_\_newindex}, \code{\_\_sub},
\code{\_\_mul}, \code{\_\_div}, \code{\_\_mod}, \code{\_\_pow},
\code{\_\_unm}, \code{\_\_lt}, \code{\_\_le}, \code{\_\_concat},
\code{\_\_len}, \code{\_\_pairs}, \code{\_\_metatable}, or weak
tables. Metatables work on tables and structs, not on lists/strings.

% ---------------------------------------------------------------------------
\chapter{Modules and Dynamic Code}
\label{cap:modulos}

\section{include}
\code{include("arquivo");} inserts the file contents at the directive
point before compilation. The path is relative to the working directory;
each file is included once per execution, which also breaks cycles. The
directive requires parentheses and a string literal; the old \code{\#}
form is not accepted.

\begin{lstlisting}
include("scripts/libs/showtext.xlb");
showtext("Ola");
\end{lstlisting}

\section{.xlb Libraries}
\code{.xlb} files are pure-language code libraries. \code{include}
embeds the content; functions and constants become available in the
script. Canonical example: \file{scripts/libs/showtext.xlb} and
\file{scripts/script\_26\_biblioteca\_xlb.isi}.

\section{require}
\code{require("arquivo")} loads an entire script as a module. The body
runs once; the file's functions and variables become globals; the body's
\code{return} value (or \code{0}) is returned and cached. \code{require}
is resolved at compile time and is not a first-class value.

\begin{lstlisting}
r = require("scripts/libs/util.xlb");
println(util_dobro(21));
r2 = require("scripts/libs/util.xlb");  // cached
\end{lstlisting}

A cycle \code{A -> B -> A} is a compilation error. A missing file or
invalid syntax is also a compilation error (not catchable by
\code{pcall}).

\section{load and dofile}
\code{load(codigo[, nome])} compiles a string as a chunk and returns a
function. In turn, \code{dofile(caminho)} reads, compiles, and executes a
file in isolation, returning the \code{return} value (or \code{0}).

\begin{lstlisting}
f = load("return 1 + 2;");
println(f());                         // 3
r = dofile("tests/fixtures/l14_chunk.txt");
\end{lstlisting}

\code{eval(codigo)} is pure sugar for \code{load(codigo)()}: it compiles
and runs the chunk immediately, returning the result. Both require
\code{allow\_modules} and \code{allow\_fileio} in the sandbox.
Loaded chunks do not capture the caller's locals; only locals from the
same lexical compilation scope. Do not use \code{load}/\code{dofile}
concurrently with threads.

\section{Native .mod Modules}
\code{import("arquivo.mod")} loads a JSON manifest that points to a
shared library (\code{.so}/\code{.dylib}). Native functions and variables
enter the global namespace.

\begin{lstlisting}
import("scripts/modules/motor_controller.mod");
println(MotorSpeedLimit);
MotorSpeedLimit = 5000;
r = Motor_Start(1500);
res = Motor_Read();
\end{lstlisting}

Accepted FFI types: signed and unsigned integers, \code{float32},
\code{float64}, \code{bool}, \code{string}, and \code{void}; at most 3
parameters. Variables with \code{readonly: true} cannot be written. See
\file{docs/linguagem.md} for the complete manifest format and module
lifecycle.

% ---------------------------------------------------------------------------
\chapter{Error Handling}
\label{cap:erros}

\section{Compilation and Runtime Errors}
Compilation errors prevent execution and include the line. Runtime errors
can be caught with \code{pcall}. \code{error(msg)} raises an error;
\code{assert(cond[, msg])} raises if the condition is false.

\begin{lstlisting}
fn valida(x) {
    assert(x > 0, "x must be positive");
    return x;
}

r = pcall(valida, -1);
println(r[0]);   // 0
println(r[1]);   // x must be positive (line ...)
\end{lstlisting}

\section{Traceback}
When an error is not caught, the output shows the message and frames:
\begin{lstlisting}[language=bash]
ERRO EM TEMPO DE EXECUCAO: Divisao por zero (linha 7)
  em quebra (linha 10)
  em main (linha 14)
\end{lstlisting}

\code{pcall} returns the message already including the traceback in
\code{r[1]}.

\section{No try/catch}
v1.0 has no \code{try}/\code{catch}/\code{finally} or \code{throw}.
Use \code{pcall} to catch errors and \code{error} to raise them. For
resource cleanup, close files explicitly or rely on automatic closing at
the end of the script/REPL.

\section{Sandbox}
The embedding API can limit instructions and disable I/O and modules. In
the runner/REPL, the REPL growth limits are configurable through the
environment. \code{load}/\code{dofile}/\code{include}/
\code{require}/\code{import(...)} respect the sandbox permissions.

% ---------------------------------------------------------------------------
\chapter{Threads and Coroutines}
\label{cap:concorrencia}

\section{System Threads}
A function declared with \code{thread} runs on a system thread. It does
not receive parameters.

\begin{lstlisting}
global n;
n = 0;

thread contador() {
    while (1) {
        n += 1;
        delay(100);   // milliseconds
    }
}

start contador();
delay(350);
stop contador();
println(n);           // approximately 3
\end{lstlisting}

\code{start}/\code{stop}/\code{pause}/\code{resume} control the thread.
Globals are shared without a lock; avoid concurrent writes to the same
variable. Graphics and input can only be used on the main thread. A
successful script waits for threads to finish; on error, they are
terminated.

\section{Cooperative Coroutines}
Coroutines use their own fixed-stack VM and are independent of system
threads.

\begin{lstlisting}
fn conta() {
    i = 0;
    while (i < 3) {
        coyield(i);
        i += 1;
    }
    return "fim";
}

co = cocreate(conta);
println(costatus(co));    // new

r = coresume(co);
println(r[0], r[1]);      // 1 0
r = coresume(co);
println(r[0], r[1]);      // 1 1
r = coresume(co);
println(r[0], r[1]);      // 1 2
r = coresume(co);
println(r[0], r[1]);      // 1 fim
println(costatus(co));    // dead
\end{lstlisting}

\code{coresume} returns \code{[ok, value]}. \code{coyield(v)} suspends
and returns \code{v}; outside a coroutine it is an error. \code{costatus}
returns \code{"novo"}, \code{"rodando"}, \code{"suspenso"}, or
\code{"morto"}. Each coroutine uses a fixed stack; create few in embedded
environments. There is no multi-value \code{coyield} and no
\code{coyield} inside \code{pcall}.

% ---------------------------------------------------------------------------
\chapter{Standard Library}
\label{cap:stdlib}

The library consists of global builtins. There are no \code{math.*},
\code{string.*}, \code{table.*}, etc.\ namespaces in v1.0. Appendix
~\ref{ap:builtins} provides a quick reference; the most used functions
are below.

\section{Console and Conversion}
\begin{lstlisting}
println("texto", 42, nil);     // text 42 nil
printf("%s=%d\n", "x", 7);     // x=7
nome = input("nome? ");        // reads a line from stdin
println(type(nome));           // string
println(tostring(42));         // 42
println(tonumber("ff", 16));   // 255
\end{lstlisting}

\code{println} accepts several arguments and prints them one by one,
separated by spaces. \code{printf} does not add a newline; use
\code{\textbackslash n}.

\section{Strings, Lists, and Tables}
\begin{lstlisting}
s = "  Ola  ";
println(upper(s));       // "  OLA  "
println(trim(s));        // "Ola"
println(sub("abcdef", 1, 3)); // bcd (inclusive end)
println(byte("A"));      // 65
println(char(65, 66));   // AB
println(format("%02d", 7));   // error: no flags/width

v = [3, 1, 2];
sort(v);                 // [1, 2, 3]
println(concat(v, "-")); // 1-2-3

t = {a = 1};
println(has(t, "a"));    // 1
println(keys(t));        // sorted list
\end{lstlisting}

\section{Mathematics and Complex Numbers}
\begin{lstlisting}
println(abs(-3));        // 3
println(sqrt(9));        // 3
println(pow(2, 10));     // 1024
println(floor(3.7));     // 3
println(ceil(3.1));      // 4
println(round(3.5));     // 4
println(min(4, 2, 9));   // 2
println(max(4, 2, 9));   // 9
println(random());       // [0, 1)
println(random_int(1, 6));

z = 3 + 4j;
println(abs(z));         // 5
println(real(z));        // 3
println(imag(z));        // 4
println(conj(z));        // 3-4j
\end{lstlisting}

\subsection*{C99 math.h}
In addition to the basic functions, the library exposes a C99 subset of
\file{math.h} for \code{double} (complex numbers only where support
already existed).

\begin{lstlisting}
println(cbrt(27));          // 3
println(hypot(3, 4));       // 5
println(ldexp(0.5, 4));     // 8
println(fma(2, 3, 4));      // 10
a, b = frexp(8);            // 0.5 and 4
println(a);                 // 0.5
println(b);                 // 4
println(fpclassify(NAN));   // 0 (FP_NAN)
println(isfinite(INF));     // 0
println(M_PI);              // ~3.14159
\end{lstlisting}

The list includes \code{acos}, \code{asin}, \code{atan}, \code{acosh},
\code{asinh}, \code{atanh}, \code{cosh}, \code{sinh}, \code{tanh},
\code{exp2}, \code{expm1}, \code{log2}, \code{log1p}, \code{cbrt},
\code{erf}, \code{erfc}, \code{tgamma}, \code{lgamma}, \code{trunc},
\code{rint}, \code{nearbyint}, \code{fabs}, \code{hypot}, \code{fdim},
\code{fmax}, \code{fmin}, \code{copysign}, \code{nextafter},
\code{remainder}, \code{lround}, \code{llround}, \code{lrint},
\code{llrint}, \code{ilogb}, \code{logb}, \code{ldexp}, \code{scalbn},
\code{scalbln}, \code{fma}, \code{nan}, \code{fpclassify}, \code{isfinite},
\code{isinf}, \code{isnan}, \code{isnormal}, \code{signbit},
\code{isgreater}, \code{isgreaterequal}, \code{isless},
\code{islessequal}, \code{islessgreater}, and \code{isunordered}.
\code{frexp}, \code{modf}, and \code{remquo} return multiple values.

Constants are literals resolved at compile time and cannot be assigned: \code{M\_E}, \code{M\_LOG2E}, \code{M\_LOG10E}, \code{M\_LN2},
\code{M\_LN10}, \code{M\_PI}, \code{M\_PI\_2}, \code{M\_PI\_4},
\code{M\_1\_PI}, \code{M\_2\_PI}, \code{M\_2\_SQRTPI}, \code{M\_SQRT2},
\code{M\_SQRT1\_2}, \code{M\_PHI}; aliases \code{PI}, \code{E},
\code{TAU}, \code{PHI}, \code{DEG\_TO\_RAD}, \code{RAD\_TO\_DEG},
\code{INF}, \code{NAN}; limits \code{HUGE\_VAL}, \code{INFINITY},
\code{MAXFLOAT}, \code{FLT\_MAX}, \code{DBL\_MAX}, \code{FLT\_EPSILON},
\code{DBL\_EPSILON}; classification \code{FP\_NAN} (0),
\code{FP\_INFINITE} (1), \code{FP\_ZERO} (2), \code{FP\_SUBNORMAL} (3),
\code{FP\_NORMAL} (4); and \code{MATH\_ERRNO} (1),
\code{MATH\_ERREXCEPT} (2), \code{math\_errhandling} (3).

\section{Time}
\begin{lstlisting}
println(time());        // seconds since the epoch
println(clock());       // milliseconds since startup
println(os_date());     // formatted local date/time
\end{lstlisting}

\section{Files}
Handles are monotonic integers; a failing \code{fopen} returns \code{0}
(the \code{NULL} constant). Argument order follows C: data first, file
last, except \code{fprintf(f, fmt, ...)}.

\begin{lstlisting}
f = fopen("dados.txt", "w");
if (f == NULL) {
    println("nao abriu");
} else {
    fputs("linha 1\n", f);
    fprintf(f, "%s=%d\n", "x", 42);
    fclose(f);
}

r = fopen("dados.txt", "r");
println(fread(100, r));
println(fgets(r));
fclose(r);
\end{lstlisting}

\section{Sound}
\begin{lstlisting}
sound_backend("device");
h = sound_load("a.wav");
sound_play(h);
sound_set_volume(h, 0.5);
delay(1000);
sound_stop(h);
sound_unload(h);
\end{lstlisting}

Real audio output is optional: use \code{sound\_backend("device")} in the
script or \code{MYINTERP\_SOUND=device}. Without that (or without raylib),
the backend is software and deterministic. \code{sound\_backend("software")}
returns to silent mode; \code{MYINTERP\_SOUND=software} also forces
software.

% ---------------------------------------------------------------------------
\chapter{Graphics}
\label{cap:graficos}

The model is inspired by Processing. The \code{software} backend is
headless and writes \code{.ppm} images; the \code{raylib} backend opens a
real window when compiled with support.

\section{setup and draw}
If the script defines \code{setup()}, it runs before any other statement.
Use \code{setDraw(fn)} (or \code{onDraw}) to register the frame
callback.

\begin{lstlisting}
void setup() {
    size(320, 200);
    setDraw(draw);
}

void draw() {
    background(20, 20, 40);
    fill(255, 160, 0);
    circle(mouseX(), mouseY(), 20);
    text("q para sair", 5, 5);
    if (keyPressed() && key() == "q") { quit(); }
}
\end{lstlisting}

Without \code{draw}/\code{setDraw}, use the manual loop:
\begin{lstlisting}
size(80, 60);
i = 0;
while (i < 60 && show()) {
    background(0);
    fill(255, 200, 0);
    circle(10 + i, 30, 12);
    i += 1;
}
\end{lstlisting}

\section{Drawing}
\begin{lstlisting}
size(80, 60);
background(20, 20, 40);
fill(255, 128, 0);
stroke(255, 255, 255);
strokeWeight(1);
noFill(); noStroke();
rect(5, 5, 30, 20);       // top-left corner
circle(60, 15, 20);       // center, diameter
ellipse(40, 45, 30, 14);  // center, axes
line(0, 59, 79, 0);
point(79, 59);
textSize(1);
font("8x8");              // "8x8" (default) or "5x7"
fill(255);
text("Isiliel", 5, 40);   // "\n" breaks the line
show();
save_screen("out.ppm");   // .ppm; .png with raylib
println(width(), height());
println(frameCount());
\end{lstlisting}

\section{Input and Images}
\begin{lstlisting}
println(mouseX(), mouseY());
println(mousePressed());
println(keyPressed());
println(key());

img = loadImage("foto.ppm");
image(img, 10, 10);
image(img, 10, 10, 64, 40);
println(imageWidth(img), imageHeight(img));
unloadImage(img);
\end{lstlisting}

Multiple keys in the same frame are coalesced: \code{key()} returns only
the last one. ESC and the X button do not close the window; use
\code{quit()} or Ctrl+C. Drawing without \code{size()}, or calling
\code{backend()} after \code{size()}, raises an error with the line
number.

% ---------------------------------------------------------------------------
\chapter{Embedding the Interpreter in C++}
\label{cap:embed}

The public API is in \file{interp.h}. Minimal example:

\begin{lstlisting}[language=C++]
#include "interp.h"

int main() {
    InterpOptions opts;
    opts.capture = true;          // stdout -> result.output
    opts.step_limit = 100000;     // instruction budget
    opts.allow_fileio = false;    // sandbox
    opts.allow_modules = false;

    InterpResult r = interp_run("println(6 * 7);", opts);
    if (!r.ok) {
        // r.error contains the message
        return 1;
    }
    // r.output contains "42\n"
    return 0;
}
\end{lstlisting}

\section{Options}
\begin{center}
\begin{tabular}{L{4cm}L{10cm}}
\toprule
Field & Meaning \\
\midrule
\code{verbose} & prints execution banners \\
\code{capture} & redirects stdout to \code{InterpResult.output} \\
\code{step\_limit} & maximum VM instructions (0 = no limit) \\
\code{allow\_fileio} & enables \code{fopen}, \code{include}, \code{require} \\
\code{allow\_modules} & enables \code{import(...)} and native modules \\
\bottomrule
\end{tabular}
\end{center}

\section{API Functions}
\begin{itemize}
  \item \code{interp\_run(source, opts)}: compiles and runs; always
        restores threads, files, screen, modules, and sandbox on exit.
  \item \code{interp\_register\_fn(nome, arity, handler)}: registers a
        host function as a builtin. Call before running.
  \item \code{interp\_set\_output(std::ostream* os)}: redirects output;
        \code{nullptr} restores \code{std::cout}.
  \item \code{interp\_reset()}: clears symbols, bytecode, threads,
        files, screen, and modules for the next single-state run.
  \item \code{run\_pipeline}, \code{run\_file}, \code{disasm\_file}:
        low-level variants used by the CLI.
\end{itemize}

\section{Builtin Signature}
A registered function receives an array of \code{Value}, the argument
count, and a pointer to the result. See \file{vm.h} for
\code{BuiltinHandler} and \file{defs.h} for \code{Value}. Errors should
be thrown as \code{std::runtime\_error}; the VM appends the line.

% ---------------------------------------------------------------------------
\chapter{Tools and Debugging}
\label{cap:ferramentas}

\section{Disassembler}
\begin{lstlisting}[language=bash]
./build/isirun --disasm script.isi
\end{lstlisting}

Shows the bytecode with opcode names and operands. Useful for
understanding optimizations and debugging the VM.

\section{Tests and Sanitizers}
\begin{lstlisting}[language=bash]
make test
make asan
make fuzz FUZZ_ITERS=2000
make bench
\end{lstlisting}

The suite covers scripts, expected errors, the REPL, the disassembler,
and the API. The fuzzer performs deterministic mutation of scripts;
ASan/UBSan check memory and undefined behavior. The benchmark compares
against Lua 5.4 and writes results to \file{bench/comparacao.html}.

\section{Project Structure}
\begin{center}
\begin{tabular}{L{4cm}L{10cm}}
\toprule
File & Role \\
\midrule
\file{lexer.cpp/h} & lexical analysis \\
\file{compiler.cpp/h} & parser and bytecode generation \\
\file{vm.cpp/h} & stack VM and builtins \\
\file{symbols.cpp/h} & symbol table and scopes \\
\file{interp.cpp/h} & pipeline and embedding API \\
\file{run.cpp} & CLI and REPL \\
\file{graphics/} & graphics HAL and backends \\
\file{scripts/} & runnable examples \\
\file{tests/} & automated suite \\
\bottomrule
\end{tabular}
\end{center}

\section{Debugging in VS Code}
The project includes an \code{F5} configuration to build and debug a
script. The debug build uses \code{-O0 -g3}; for performance use
\code{make release}.

% ---------------------------------------------------------------------------
\chapter{Known Limitations and Differences}
\label{cap:limitacoes}

This section summarizes what does \emph{not} exist or differs from the
project documentation in v1.0.

\section{Language}
\begin{itemize}
  \item \code{true}/\code{false} are 1/0; \code{type(true)} returns
        \code{number}. The internal \code{VAL\_BOOL} type exists, but it is
        not produced by literals/comparisons in v1.0.
  \item \code{\&\&}/\code{||} do not short-circuit and return 1/0.
        \code{!} and \code{\textasciitilde} are bitwise negation; there
        is no logical negation.
  \item There is no \code{local}, \code{let}, or \code{const};
        assignment inside a function is implicitly local.
  \item Only C-style \code{for} exists. There is no \code{for k,v in ...}
        or native iterators; \code{pairs}/\code{ipairs} return lists.
  \item There is no \code{and}/\code{or}/\code{not}, \code{\#},
        \code{..}, or \code{//} operator. Use \code{floordiv} and
        \code{+}.
  \item There is no IIFE or postfix call/indexing on arbitrary
        expressions: \code{f()[0]} and \code{fn(x)\{...\}(1)} are not
        supported. Use temporaries.
  \item Named functions and methods do not capture outer locals; only
        anonymous functions do. There is no forward capture.
  \item Table literals accept only identifier keys; use
        \code{t[chave] = v} for dynamic keys.
  \item There are no library namespaces (\code{math.*}, \code{string.*},
        \code{table.*}) or first-class builtins.
  \item There is no \code{try}/\code{catch}/\code{finally}/\code{throw}.
        Use \code{pcall}.
\end{itemize}

\section{Metatables and Library}
\begin{itemize}
  \item Missing: \code{\_\_newindex}, \code{\_\_sub}, \code{\_\_mul},
        \code{\_\_div}, \code{\_\_mod}, \code{\_\_pow}, \code{\_\_unm},
        \code{\_\_lt}, \code{\_\_le}, \code{\_\_concat}, \code{\_\_len},
        \code{\_\_pairs}, \code{\_\_metatable}, and weak tables.
  \item Metatables do not work on lists or strings.
  \item \code{match} returns \code{""} on failure (not \code{nil});
        \code{gmatch} returns a list (not an iterator); \code{gsub} only
        accepts textual replacement.
  \item \code{printf}/\code{format} do not accept width, precision, or
        flags.
  \item \code{require} is resolved at compile time and is not a
        first-class value; there is no \code{package.path} or
        \code{package.loaded}.
  \item There is no \code{debug.*} or stack introspection.
\end{itemize}

\section{Concurrency and Memory}
\begin{itemize}
  \item The cycle collector is not concurrent; cycles in paused
        thread/coroutine stacks may be retained.
  \item \code{load}/\code{dofile} bytecode accumulates until the end of
        execution; there is no per-chunk release.
  \item Threads share globals without a lock; concurrent writes are the
        user's responsibility.
  \item Coroutines have a fixed stack and single-value \code{coyield};
        there is no \code{coyield} inside \code{pcall}.
\end{itemize}

\section{Roadmap}
The only structural item still pending in the documents is E18 (register
VM), which is a performance optimization. The functional gaps above are
candidates for future phases: real booleans, short-circuiting, block
scope, named closures, postfix/IIFE, \code{for-in}, namespaces, complete
metatables, and error handling with \code{try/catch}.

% ===========================================================================
\appendix
% ===========================================================================

\chapter{Grammar Summary}
\label{ap:gramatica}

The complete grammar is in \file{docs/gramatica.ebnf}. Summary:

\begin{lstlisting}[language=]
program        = { statement } EOF ;
block          = "{", { statement }, "}" ;

statement      = block | if_statement | while_statement
               | for_statement | repeat_statement
               | "return", [ expression_list ], ";"
               | "break", ";" | "continue", ";"
               | "global", identifier, { ",", identifier }, ";"
               | "goto", identifier, ";"
               | "::", identifier, "::"
               | type_definition | class_definition
               | thread_definition | thread_command
               | include_directive | import_directive
               | function_definition | scalar_declaration
               | expression_statement ;

if_statement   = "if", "(", expression, ")", statement,
                 { "elseif", "(", expression, ")", statement },
                 [ "else", statement ] ;

while_statement = "while", "(", expression, ")", statement ;

for_statement  = "for", "(", [ expression ], ";",
                 [ expression ], ";",
                 [ expression ], ")", statement ;

repeat_statement = "repeat", statement, "until",
                   [ "(" ], expression, [ ")" ], [ ";" ] ;

function_definition = [ type_prefix ], [ "fn" ], identifier,
                      "(", [ parameter_list ], ")", block ;

anonymous_function = [ "fn" ], "(", [ parameter_list ], ")", block ;

type_definition  = "type", identifier, "{", { struct_member }, "}", ";" ;
class_definition = "class", identifier, [ "extends", identifier ],
                   "{", { struct_member }, "}", [ ";" ] ;
\end{lstlisting}

Precedence (lowest to highest): assignment, \code{||},
\code{\&\&}, \code{|}, \code{\^{}}, \code{\&}, \code{== !=},
\code{< > <= >=}, \code{<< >>}, \code{+ -}, \code{* / \%}, unary,
postfix.

\chapter{Quick Builtin Reference}
\label{ap:builtins}

This list covers the most used functions. For the complete list and
arity details, see \file{docs/linguagem.md} and the code in
\file{vm.cpp}, \file{patterns.cpp}, \file{fileio.cpp},
\file{threads.cpp}, \file{sound.cpp}, and \file{graphics/}.

\begin{longtable}{@{}L{4.8cm}L{10.5cm}@{}}
\toprule
Function & Summary \\
\midrule
\endhead
\code{println(...)} & prints values with a newline \\
\code{printf(fmt, ...)} & prints formatted output, no newline \\
\code{format(fmt, ...)} & returns a formatted string \\
\code{input([prompt])} & reads a line from stdin \\
\code{type(x)} & type name \\
\code{str(x)}, \code{tostring(x)} & textual representation \\
\code{tonumber(x[, base])} & converts to a number or \code{nil} \\
\code{assert(cond[, msg])} & error if the condition is false \\
\code{error(msg)} & raises an error \\
\code{collectgarbage()} & runs the cycle collector \\
\midrule
\code{len(x)} & length of string/list/table \\
\code{append(lista, x)} & appends to the end of a 1D list \\
\code{sort(v)} & sorts numbers (in place) \\
\code{concat(v[, sep])} & joins elements into text \\
\code{upper(s)}, \code{lower(s)} & uppercase/lowercase \\
\code{trim(s)} & removes leading/trailing spaces/tabs \\
\code{rep(s, n)} & repeats a string \\
\code{reverse(s)} & reverses a string \\
\code{find(s, sub)} & 0-based index or -1 \\
\code{sub(s, i[, j])} & 0-based substring, inclusive end \\
\code{byte(s[, i])} & byte code \\
\code{char(b1, ...)} & string from bytes \\
\code{match(s, pat[, init])} & first match/capture \\
\code{gmatch(s, pat)} & list of matches \\
\code{gsub(s, pat, repl)} & textual substitution \\
\code{utf8\_len(s)} & code points \\
\code{utf8\_codepoint(s[, i])} & code point or -1 \\
\code{utf8\_char(cp)} & encodes a code point \\
\code{String(...)} & text constructor (float: 2 decimals; int: base) \\
\code{charAt}, \code{compareTo}, \code{c\_str} & Arduino String API \\
\code{endsWith}, \code{equals}, \code{equalsIgnoreCase} & text comparisons \\
\code{getBytes}, \code{toCharArray} & byte list or buffer \\
\code{indexOf}, \code{lastIndexOf}, \code{length} & search and length \\
\code{remove(s,i[,n])}, \code{replace}, \code{setCharAt} & immutable editing \\
\code{reserve}, \code{startsWith}, \code{substring} & Arduino utilities \\
\code{toDouble}, \code{toFloat}, \code{toInt} & text conversion \\
\code{toLowerCase}, \code{toUpperCase} & upper/lower case \\
\midrule
\code{has(t, k)} & 1 if the key exists \\
\code{remove(t, k)} & removes a key; 1 if it existed \\
\code{keys(t)} & sorted list of keys \\
\code{rawget(t, k)}, \code{rawset(t, k, v)} & access without metamethods \\
\code{next(t[, k])} & next pair or empty list \\
\code{pairs(t)} & list of pairs \\
\code{ipairs(t)} & values of \code{t[0..]} \\
\code{select(n, ...)} & selects arguments \\
\code{unpack(v[, i[, j]])} & returns elements \\
\midrule
\code{abs}, \code{sqrt}, \code{pow} & mathematical functions \\
\code{sin}, \code{cos}, \code{tan} & trigonometry (complex numbers) \\
\code{atan2(y, x)} & arc tangent \\
\code{log}, \code{log10}, \code{exp} & logarithms/exponential \\
\code{floor}, \code{ceil}, \code{round} & rounding \\
\code{floordiv(a, b)} & integer division \\
\code{fmod(a, b)} & remainder with dividend sign \\
\code{min(...)}, \code{max(...)} & minimum/maximum \\
\code{random()}, \code{random\_int(a,b)} & random numbers \\
\code{conj}, \code{real}, \code{imag} & complex numbers \\
\code{rect(re,im)}, \code{polar(m,ang)} & complex construction \\
\code{acos}, \code{asin}, \code{atan} & inverse trigonometry \\
\code{cosh}, \code{sinh}, \code{tanh} & hyperbolic trigonometry \\
\code{acosh}, \code{asinh}, \code{atanh} & inverse hyperbolic \\
\code{cbrt}, \code{erf}, \code{erfc} & cube root and error \\
\code{tgamma}, \code{lgamma} & gamma functions \\
\code{exp2}, \code{expm1}, \code{log2}, \code{log1p} & extra exp/log \\
\code{trunc}, \code{rint}, \code{nearbyint} & C99 rounding \\
\code{hypot}, \code{fdim}, \code{fmax}, \code{fmin} & C99 binary \\
\code{copysign}, \code{nextafter}, \code{remainder} & C99 binary \\
\code{lround}, \code{llround}, \code{lrint}, \code{llrint} & integer rounding \\
\code{ilogb}, \code{logb}, \code{ldexp}, \code{scalbn}, \code{scalbln} & exponent/scale \\
\code{fma}, \code{nan} & C99 \\
\code{frexp}, \code{modf}, \code{remquo} & multiple return \\
\code{fpclassify}, \code{isnormal}, \code{signbit} & classification \\
\code{isgreater}, \code{isless}, \code{isunordered} & C99 comparison \\
\code{PI}, \code{E}, \code{TAU}, \code{M\_PI}, \code{M\_*} & mathematical constants \\
\code{FP\_*}, \code{math\_errhandling} & classification/environment \\
\midrule
\code{time()}, \code{clock()} & clocks \\
\code{os\_time()}, \code{os\_clock()} & aliases \\
\code{os\_date([fmt])} & local date/time \\
\midrule
\code{mat\_index(M,i,j)} & 1-based linear index \\
\code{rows(M)}, \code{cols(M)} & dimensions \\
\code{identity(n)}, \code{zeros(n)} & special matrices \\
\code{det}, \code{trace}, \code{inv}, \code{transpose} & linear algebra \\
\code{dot}, \code{cross}, \code{matmul} & products \\
\code{eigvals(M)}, \code{eigvecs(M)} & eigenvalues/eigenvectors \\
\midrule
\code{fopen}, \code{fclose} & open/close file \\
\code{fread}, \code{fwrite}, \code{fputs} & read/write \\
\code{fgets}, \code{fgetc}, \code{fputc} & line/character \\
\code{fseek}, \code{ftell}, \code{rewind} & position \\
\code{feof}, \code{ferror} & state \\
\code{fprintf(f, fmt, ...)} & formatted write \\
\midrule
\code{cocreate(f)} & creates a coroutine \\
\code{coresume(co, ...)} & resumes a coroutine \\
\code{coyield(v)} & suspends a coroutine \\
\code{costatus(co)} & coroutine state \\
\code{start/stop/pause/resume f()} & controls a thread \\
\code{delay(ms)} & sleeps for milliseconds \\
\midrule
\code{load(src[, nome])} & compiles a chunk \\
\code{dofile(path)} & executes a file \\
\code{require("path")} & loads a module (compile time) \\
\midrule
\code{setmetatable(t, mt)} & sets a metatable \\
\code{getmetatable(t)} & gets a metatable \\
\code{rawget}, \code{rawset} & direct access \\
\midrule
\code{size(w,h)}, \code{background(...)} & graphics \\
\code{fill(...)}, \code{stroke(...)} & colors \\
\code{rect}, \code{circle}, \code{ellipse}, \code{line} & primitives \\
\code{text(...)}, \code{font(...)}, \code{textSize(...)} & text \\
\code{show()}, \code{quit()} & frame/termination \\
\code{loadImage}, \code{image}, \code{unloadImage} & images \\
\code{mouseX}, \code{mouseY}, \code{keyPressed}, \code{key} & input \\
\code{sound\_load}, \code{sound\_play}, \code{sound\_stop} & sound \\
\bottomrule
\end{longtable}

\chapter{Complete Programs}
\label{ap:exemplos}

\section{FizzBuzz}
\begin{lstlisting}
for (i = 1; i <= 30; i += 1) {
    if (i % 15 == 0) {
        println("FizzBuzz");
    } elseif (i % 3 == 0) {
        println("Fizz");
    } elseif (i % 5 == 0) {
        println("Buzz");
    } else {
        println(i);
    }
}
\end{lstlisting}

\section{Iterative and Recursive Fibonacci}
\begin{lstlisting}
fib_iter(n) {
    a = 0; b = 1;
    for (i = 0; i < n; i += 1) {
        t = a + b;
        a = b;
        b = t;
    }
    return a;
}

fib_rec(n) {
    if (n < 2) { return n; }
    return fib_rec(n - 1) + fib_rec(n - 2);
}

println(fib_iter(20));   // 6765
println(fib_rec(10));    // 55
\end{lstlisting}

\section{Counter Closure}
\begin{lstlisting}
contador() {
    n = 0;
    return fn() {
        n += 1;
        return n;
    };
}

c = contador();
println(c());   // 1
println(c());   // 2
println(c());   // 3
\end{lstlisting}

\section{Table and Iteration}
\begin{lstlisting}
estoque = {
    maca = 10,
    banana = 5,
    laranja = 8
};

ks = keys(estoque);
for (i = 0; i < len(ks); i += 1) {
    k = ks[i];
    println(k, "=", estoque[k]);
}

ps = pairs(estoque);
for (i = 0; i < len(ps); i += 1) {
    println(ps[i][0], ps[i][1]);
}
\end{lstlisting}

\section{Class with Inheritance}
\begin{lstlisting}
class Animal {
    nome;
    init(n) { nome = n; }
    falar() { return "..."; }
}

class Cachorro extends Animal {
    init(n) { super.init(n); }
    falar() { return "au au"; }
}

a = new Animal("bicho");
d = new Cachorro("rex");
println(a.falar());   // ...
println(d.falar());   // au au
println(d.nome);      // rex
\end{lstlisting}

\section{Metatable for Addition}
\begin{lstlisting}
mt = {__add = fn(a, b) { return a.x + b.x; }};
a = setmetatable({x = 2}, mt);
b = setmetatable({x = 3}, mt);
println(a + b);   // 5
\end{lstlisting}

\section{Coroutine}
\begin{lstlisting}
fn produtor() {
    for (i = 1; i <= 3; i += 1) {
        coyield(i * 10);
    }
    return "fim";
}

co = cocreate(produtor);
while (1) {
    r = coresume(co);
    println(r[1]);
    if (r[1] == "fim") { break; }
}
\end{lstlisting}

\section{Simple Graphics}
\begin{lstlisting}
void setup() {
    size(160, 120);
    setDraw(quadro);
}

void quadro() {
    background(15, 15, 35);
    fill(255, 180, 0);
    circle(80, 60, 40);
    if (keyPressed() && key() == "q") { quit(); }
}
\end{lstlisting}

\chapter{Glossary}
\label{ap:glossario}

\begin{description}[leftmargin=3cm,style=nextline]
  \item[Array] 0-based list/matrix, represented by \code{VAL\_ARRAY}.
  \item[Builtin] Native function registered in the VM and called by name.
  \item[Bytecode] Instructions produced by the compiler and executed by the VM.
  \item[Closure] Anonymous function that captures local variables from the enclosing scope.
  \item[Coroutine] Cooperative execution that can be suspended and resumed.
  \item[Metatable] Table that defines metamethods for another table/struct.
  \item[Metamethod] Special field such as \code{\_\_index} or \code{\_\_add}.
  \item[REPL] Interactive read-eval-print environment.
  \item[Sandbox] Restrictions on I/O, modules, and instruction count.
  \item[Struct] Type with fields and methods, defined by \code{type}.
  \item[Table] Key-value map with reference semantics.
  \item[Stack VM] Virtual machine that uses a stack of values.
\end{description}

\end{document}
